\gddchapter{Testing sequence}
The prototype can be played with both voice and keyboard and the navigation modes can be swapped with a button press. 

\section{Baseline A - Keyboard navigation}
At first the participant will spend 10 minutes by playing the game using keyboard navigation.
This allows the player to anchor in the typical experience and serves as a comparative point in the A/B testing.

\subsection{Observation log}
\begin{table}[h]
\centering
\begin{tabular}{|p{5cm}|p{10cm}|}
\hline
\textbf{Field} & \textbf{Response} \\ \hline
\textbf{Physical comfort} How is the participant feeling & The participant seems to be comfortable \\[1.5cm] \hline
\textbf{Immersion markers} (flow \& frustration) & The participant seemed to be deeply engaged with the eperience in both versions of the game\\[1.5cm] \hline
\end{tabular}
\end{table}

\section{Novel version B - Voice navigation}
After 10 minutes the researcher switches the navigation to voice input mode by using the keyboard. The prototype is not restarted,
instead the player is instructed to continue his playthrough using voice. The decision not to restart is placed to reduce any frustration that the player
might feel if they have to restart and repeat the actions that they just did.

\subsection{Observation log}
\begin{table}[H]
\centering
\begin{tabular}{|p{5cm}|p{10cm}|}
\hline
\textbf{Field} & \textbf{Response} \\ \hline
\textbf{Physical comfort} How is the participant feeling & Again, the participant seems to be comfortable\\[1.5cm] \hline
\textbf{Natural language observation} (is the player using long sentences or simple commands? What type of language do they use? Are they speaking first or 3rd person?) & The player is using simple commands, targeting 2nd person (eg. go to the lounge, pick up the crowbar)\\[1.5cm] \hline
\textbf{Linguistic guessing} (is the particpant trying to guess the possible commands by speaking to the game?) & The participant seems to know the pre-exisitng limits from the keyboard version. She's not trying to do things that are beyond the limitations of the keyboard experience.\\[1.5cm] \hline
\textbf{Processing time user reaction} (Does the wait time break the suspension of disbelief)& Defenietly a little bit. It seems like wait time gets the participant out of the experience a little bit.\\[1.5cm] \hline
\textbf{Voice processing pipeline edge cases} (System edge cases eg. two vs to)& The engine didn't understand the word "lounge" very well. This being one of the game locations presented an issue\\[1.5cm] \hline
\end{tabular}
\end{table}

\subsection{Sticky moments}
The moments that stood out during testing marked with timestamps.
\begin{table}[h]
\centering
\begin{tabular}{|p{3cm}|p{10cm}|}
\hline
\textbf{Timestamp} & \textbf{Log} \\ \hline
05:00 & Start of the playtest \\ \hline 
06:00 & The player complaining about the length of the location descriptions \\ \hline 
10:45 & It's not directly obvious how to use items for the player in the keyboard version \\ \hline 
14:05 & The player complains again about reading time needed to play "I could be not reading this at all. This takes too long" \\ \hline 
16:05 & The naming convention for the stairs is not intuitive \\ \hline 
17:00 & The UX of recording is not very obvious \\ \hline 
18:30 & Bug - the 1F stairs have the same image as 0F. I defended the game a bit too much\\ \hline 
21:00 & Double UX fail entering the lounge \\ \hline 
22:45 & Pharmacy, the leg stabiliser refused to be picked up at first but it worked on the second attempt. \\ 


\end{tabular}
\end{table}

\section{Alignment check}
How much was the researcher participating in the gameplay experience? 
Is the participant related in any way to the researcher? Could that influence the result?

The researcher participated in the form of explaining the controls and defended the game a little bit when the VUI didn't trigger correctly for Lounge. This observation was carried out to make sure that the researcher won't defend the game moving forward.
